\chapter{热传递的基本原理}

\begin{problem}
傅立叶定律说明了什么问题？其数学表达式是怎样的？
\end{problem}

\begin{solution}
\begin{enumerate}
  \item \textbf{说明的问题}：单位时间内通过导热体单位面积上的导热量，在数值上与该面积上的温度梯度成正比，而方向与其相反。
  \item \textbf{数学表达式}：\(q=-\lambda\frac{t}{x}\) 或 \(Q=-\lambda S\frac{t}{x}\)。
\end{enumerate}
\end{solution}

\begin{problem}
多层圆简壁的导热热阻如何计算？每米管长的导热量及接触面上的温度如何计算？
\end{problem}

\begin{solution}
\begin{enumerate}
  \item \textbf{多层圆筒壁的导热热阻}：单位管长上的总导热热阻为各层圆筒壁热阻之和，计算公式为
  \begin{equation*}
    \sum_{i=1}^{n}\frac{1}{2\pi\lambda_{i}}\ln\frac{r_{i+1}}{r_{i}}
  \end{equation*}
  \item \textbf{每米管长的导热量}：利用总温差除以总热阻得出，即
  \begin{equation*}
    q_{l}=\frac{t_{w1}-t_{w(n+1)}}{\sum_{i=1}^{n}\frac{1}{2\pi\lambda_{i}}\ln\frac{r_{i+1}}{r_{i}}}
  \end{equation*}
  \item \textbf{接触面上的温度}：在求出每米管长导热量 \(q_l\) 的基础上，利用相邻两层之间的热阻关系逐层推算得出。
\end{enumerate}
\end{solution}

\begin{problem}
牛顿冷却定律解决了什么问题？表面传热系数的单位与热导率有什么不同？
\end{problem}

\begin{solution}
\begin{enumerate}
  \item \textbf{解决的问题}：牛顿冷却定律给出了对流换热量的计算方法，将复杂的对流换热过程及其诸多影响因素归并至"对流传热系数"中，从而使求解对流换热问题变成了求解表面传热系数的问题。
  \item \textbf{单位区别}：表面传热系数（\(\alpha\)）的单位是 \unit{W/(m^2 \cdot \degreeCelsius)}，而热导率（\(\lambda\)）的单位是 \unit{W/(m \cdot \degreeCelsius)}。
\end{enumerate}
\end{solution}

\begin{problem}
影响对流换热的主要因素有哪些？
\end{problem}

\begin{solution}
主要包括以下五个方面：
\begin{enumerate}
  \item \textbf{流动起因}：自然对流或强制对流。
  \item \textbf{流动状态}：层流或紊流。
  \item \textbf{流体有无相变}：单相换热或相变换热（如沸腾、凝结等）。
  \item \textbf{换热表面的几何因素}：如内部流动（管内）、外部绕流（外掠平板或管束）等。
  \item \textbf{流体的热物理性质}：包括热导率、比热容、运动/动力粘度、密度和体胀系数等。
\end{enumerate}
\end{solution}

\begin{problem}
管内强制对流换热与哪些准则有关？各准则说明了什么？它们的数学表达式如何？
\end{problem}

\begin{solution}
主要与以下无量纲准则数有关：
\begin{enumerate}
  \item \textbf{努塞尔数（\(Nu\)）}：包含表面传热系数，反映了对流换热的强弱程度。表达式为 \(Nu=\frac{\alpha d}{\lambda}\)。
  \item \textbf{普朗特数（\(Pr\)）}：反映流体物理性质对换热的影响。表达式为 \(Pr=\frac{\mu c}{\lambda}\)。
  \item \textbf{雷诺数（\(Re\)）}：反映流体的流动状态（层流或紊流），强制对流换热经验公式中通常包含此项。
\end{enumerate}
\end{solution}

\begin{problem}
吸收率、反射率和透射率的含义是什么？黑体的吸收率如何？实际物体的吸收率与黑度之间有什么样的关系？
\end{problem}

\begin{solution}
\begin{enumerate}
  \item \textbf{含义}：外界投入到物体表面上的总辐射能量中，被物体所吸收、反射和穿透的能量所占的比例，分别称为物体的吸收率（\(A\)）、反射率（\(R\)）和透射率（\(D\)）。
  \item \textbf{黑体的吸收率}：吸收率 \(A=1\)。
  \item \textbf{吸收率与黑度的关系}：根据基尔霍夫定律，在热平衡条件下，实际物体的吸收率在数值上恒等于同温度下该物体的黑度（即 \(A=\epsilon\)）。
\end{enumerate}
\end{solution}

\begin{problem}
物体的本身辐射、投射辐射、吸收辐射及有效辐射的含义是什么？其间的关系如何？
\end{problem}

\begin{solution}
\begin{enumerate}
  \item \textbf{本身辐射（\(E\)）}：单位时间内，物体的单位表面积向外发射的辐射能。
  \item \textbf{投射辐射（\(E_e\)）}：被外界投入到一物体表面上的总辐射能量。
  \item \textbf{吸收辐射（\(E_a\)）}：投射辐射中被物体所吸收的能量部分。
  \item \textbf{有效辐射（\(E_{ef}\)）}：本身辐射 \(E\) 和反射辐射 \(E_r\) 的总和。
  \item \textbf{关系}：有效辐射等于本身辐射加上被反射的投射辐射部分。对于不透热体，其关系可表示为 \(E_{ef}=E+(1-A)E_e\)。
\end{enumerate}
\end{solution}

\begin{problem}
何谓传热过程？通过两层圆筒壁的传热过程包括哪几个环节？其传热热阻如何计算？
\end{problem}

\begin{solution}
\begin{enumerate}
  \item \textbf{传热过程}：热量从温度较高的流体经过固体壁传递给另一侧温度较低流体的过程，称为总传热过程（简称传热过程）。
  \item \textbf{环节}：通过两层圆筒壁的传热过程包括三个主要环节：热流体到内侧壁面的对流换热、两层固体壁的导热、外侧壁面与冷流体的对流换热。
  \item \textbf{传热热阻计算}：总传热热阻等于内部对流热阻、两层圆柱面导热热阻与外部对流热阻之和。计算公式为
  \begin{equation*}
    R_{\text{总}} = \frac{1}{\alpha_{1}\pi d_{1}l} + \frac{1}{2\pi\lambda_{1}l}\ln\frac{d_{2}}{d_{1}} + \frac{1}{2\pi\lambda_{2}l}\ln\frac{d_{3}}{d_{2}} + \frac{1}{\alpha_{2}\pi d_{3}l}
  \end{equation*}
\end{enumerate}
\end{solution}

\begin{problem}
强化传热与削弱传热的措施有哪些？
\end{problem}

\begin{solution}
\textbf{强化传热措施}：
\begin{enumerate}
  \item \textbf{增大传热系数}：降低传热过程的总热阻。手段分为无源技术（如表面涂层、粗糙表面、扩展表面、加扰流元件等）和有源技术（如机械搅拌、使换热表面或流体振动等）。
  \item \textbf{增大传热温差}：通过改变受热面的布置方式（如采用逆流布置）来实现。
  \item \textbf{增大换热面积}：采用小直径管束或在壁面上加肋等方法。
\end{enumerate}

\textbf{削弱传热措施（隔热保温技术）}：
\begin{enumerate}
  \item 主要通过控制导热过程来减少热量损失。
  \item 对高于环境温度的热力设备，多采用无机绝热材料。
  \item 对低于环境温度的设备，可选择一般性绝热材料、抽真空粉末颗粒绝热材料或多层真空绝热材料。
\end{enumerate}
\end{solution}
